\begin{description}
\item[Q1. Why not simply use simulation to find equivalent gates?] For finding candidates, simulation is the right
tool and remains far better than any embedding (\cref{sec:matching}). Retrieval is a \emph{probe}: if a representation
claims to encode function, equivalent gates should be close in it. The probe also matters downstream, because
embeddings reused for synthesis, verification or design-space exploration inherit this property.

\item[Q2. Is identity injection realistic?] Not as a flow, and we do not use it as one. It is a stress test that
controls the amount of redundancy; all claims about realistic rewrites rest on the synthesis and held-out families,
which are reported separately everywhere.

\item[Q3. The audited model is small. Do the findings transfer to DeepGate?] The bidirectional encoder adds
DeepGate's attention and reverse pass and shows the same pattern (\cref{sec:generality}), and \cref{tab:general}
applies the unchanged harness to the released DeepGate2 and DeepGate3 checkpoints, which show it too. We make no
claim about published models beyond the quantities that table reports.

\item[Q4. Does $\rho$ exaggerate the problem?] $\rho$ is reported together with the absolute inconsistency, the
noise floor and $\Delta$MAE. Under synthesis, the absolute inconsistency is \IncSynOverNoiseBase\ times the noise floor,
which does not involve the model's error.

\item[Q5. Are the improvements of \ours\ significant?] At design level, \ours\ improves retrieval on 7 of 7 test
designs against every other regime, with bootstrap intervals that exclude zero (\cref{tab:app_stats}). Its accuracy
cost is not significant at design level.

\item[Q6. Was the test set used for development?] An exploratory version of the study was evaluated on the same
test designs; we disclose this in the main paper. The final protocol, the selection rule and all reported models
were fixed and retrained afterwards, and the ISCAS suite was never used before the final evaluation.

\item[Q7. Why does H2 break every model?] H2 replaces each affected AND by XORs whose inputs reconverge. A likely reason is that signal
probabilities of reconvergent logic depend on correlations between fanins, which a message-passing encoder trained on
AND-dominated cones never had to represent. Consistent with this, the failure appears for every regime, including those
trained on certified variants, on both suites and with every encoder.

\item[Q8. Why are complemented pairs excluded from \ours\ and from retrieval?] They are \ComplPct\% of pairs. A single
embedding cannot be close both to $f$ and to $\neg f$ without a polarity-aware geometry; we leave such a geometry
to future work and keep retrieval and training on the same pair set.

\item[Q9. Are labels correct?] Signal probabilities are estimated from 4{,}096 uniform patterns at the circuit inputs
(standard error at most 0.008). Equivalences are proven by SAT, not estimated.

\item[Q10. Does the accuracy cost of \ours\ make it impractical?] With $\lambda=0.05$ it keeps baseline-level
accuracy at almost the same retrieval (\cref{tab:ablation}); the selected $\lambda$ follows a rule fixed in advance.
The trade-off can be set by the user.

\item[Q11. Why is SAT sweeping slower with embeddings?] On designs of at most 1{,}020 gates, refuted calls take
microseconds and the embedding forward pass dominates. The probe shows that the ordering signal exists; turning it
into a speed-up requires larger designs, hard refutations and amortized or incremental inference.

\item[Q12. Could the structural-hash definition of ``restructured'' miss changes?] Two cones with equal hashes are
isomorphic up to AND commutativity, so a fanin-only encoder gives them identical states. The definition is therefore
exact for such encoders; for others the all-pairs metrics apply.

\item[Q13. DeepGate2 filters its training pairs. Does TT-sim fail only because ours are sampled at random?] No. We
retrained the objective with DeepGate2's own filters, namely equal signal probability within $0.05$, level difference
at most $5$, and truth-table distance either below $0.2$ or above $0.8$, each pair drawn inside a single cone so the
inputs agree by construction; distances are zero-normalized as in DeepGate2. Over \TTFNSeeds\ seeds the filtered variant
retrieves less than the random-pair version (\TTFTopRS\% against \TTTopRS\%), is no more accurate (MAE \TTFMAE\ against
\TTMAE) and leaves synthesis unchanged ($\rho$ \TTFRhoSyn\ against \TTRhoSyn); it lowers the inconsistency under
injection slightly ($\rho$ \TTFRhoInj\ against \TTRhoInj, redundancy gap \TTFDInj\ against \TTDInj), far from \ours\
($\rho$ \OursRhoInj, redundancy gap \OursDInj, retrieval \OursTopRS\%). Filtering makes
the pairs cleaner, but it leaves the objective pairwise and functional-similarity-based, which is exactly the property
the audit finds insufficient.

\item[Q14. If simulation-matched pairs train \ours\ as well as proven pairs, why prove anything?] Training and
measurement need different things. Training tolerates a small fraction of wrong pairs: with pairs matched on
4{,}096-bit or 64-bit signatures, \ours\ reaches \SigFourKTop\% and \SigSixtyFourTop\% retrieval against
\OursTopRS\% with proven pairs. Measurement does not: \SigFalsePct\% of the gates with a signature match have no
true equivalent, and such a gate would be scored as a miss of a pair that does not exist. Proofs make every pair in
the audit a true positive.

\item[Q15. Is \ours\ just a contrastive head on top of an unchanged encoder?] No. Training only the projection head
with the same loss and pairs on a frozen encoder gives \ProbeBaseTop\% (baseline encoder) and \ProbeAugTop\%
(augmented encoder), against \OursTopRS\% for \ours; \EncShare\% of the gain over the baseline requires training the
encoder.

\item[Q16. Was the confirmatory design chosen after seeing results?] hyp was not used at any stage of development.
We built its circuits and evaluated the 15 final models only after every other result was fixed, with the same
code and no retraining; its main findings replicate there, and two secondary ones do not (\cref{sec:generality}).

\item[Q17. MixGate aligns SAT-proven pairs with an $L_1$ loss. Why InfoNCE?] Because recognizing an equivalent among
many gates needs negatives. On our proven pairs, with the weight chosen by the same rule, $L_1$ alignment reaches
\LoneTop\% retrieval ($L_1$ alignment + TT-sim: \LonettTop\%), close to augmentation, while its inconsistency matches
augmentation's; at larger weights all embeddings collapse to one point. A positives-only loss steadies the
predictions but does not make the representation discriminative.
\end{description}
